The primary objective of this laboratory exercise is to mathematically analyze and computationally simulate continuous-time signals using 
Fourier Series and Fourier Transforms. Analyzing signals in the frequency domain shows the individual sinusoidal frequency components that 
make up complex waveforms \cite{oppenheim1997}. 

The exponential Fourier Series specifically decomposes continuous-time periodic signals into an infinite sum of complex exponential 
functions \cite{riskin2003}.
Reconstructing a signal using only a finite number of exponential terms results in waveform that exhibits oscillations 
near discontinuities which is a behavior known as the Gibbs phenomenon \cite{haykin1999}.

Meanwhile for aperiodic signals the Fourier Transform extends to analyzing the signal over a continuous spectrum of frequencies instead 
of the discrete harmonics \cite{oppenheim1997}. Applying Fourier Transform to a system's impulse response gives its transfer function. The transfer function 
dictates how the system filters different input frequencies \cite{riskin2003}. Moreover, the operational boundary of this filter is defined by 3dB cut-off 
frequency that marks the exact point where the signal's power is reduced by half \cite{hsu2013}.

Several specific MATLAB functions are required to computationally implement and visualize these mathematical operations. The syms and assume 
functions establish the foundational symbolic variables and their specific mathematical constraints for calculus operations. The int() 
function calculates the exact analytical integral to find the Fourier coefficients, while sum() and subs() are combined to evaluate and add 
the truncated series terms. Finally, the fourier() command computes the continuous transfer function, and plotting tools like stem(), 
plot(), and fplot() construct the visual representations of the signal magnitudes and cut-off relationships.

